\documentclass[10pt,a4paper]{amsart}
\usepackage{amsmath,amssymb,amsthm,graphicx}
% For commutative diagrams we use Xy-pic
\usepackage[all]{xy}

% We use 2cell for 2-commutative diagrams.
\xyoption{2cell}
% Theorem environments.
%
\theoremstyle{plain}
\newtheorem{theorem}[subsection]{Theorem}
\newtheorem{proposition}[subsection]{Proposition}
\newtheorem{lemma}[subsection]{Lemma}

\theoremstyle{definition}
\newtheorem{definition}[subsection]{Definition}
\newtheorem{example}[subsection]{Example}
\newtheorem{exercise}[subsection]{Exercise}
\newtheorem{situation}[subsection]{Situation}

\theoremstyle{remark}
\newtheorem{remark}[subsection]{Remark}
\newtheorem{remarks}[subsection]{Remarks}

\numberwithin{equation}{subsection}

% Macros
%
\def\lim{\mathop{\mathrm{lim}}\nolimits}
\def\colim{\mathop{\mathrm{colim}}\nolimits}
\def\Spec{\mathop{\mathrm{Spec}}}
\def\Hom{\mathop{\mathrm{Hom}}\nolimits}
\def\Ext{\mathop{\mathrm{Ext}}\nolimits}
\def\SheafHom{\mathop{\mathcal{H}\!\mathit{om}}\nolimits}
\def\SheafExt{\mathop{\mathcal{E}\!\mathit{xt}}\nolimits}
\def\Sch{\mathit{Sch}}
\def\Mor{\mathop{\mathrm{Mor}}\nolimits}
\def\Ob{\mathop{\mathrm{Ob}}\nolimits}
\def\Sh{\mathop{\mathit{Sh}}\nolimits}
\def\NL{\mathop{N\!L}\nolimits}
\def\CH{\mathop{\mathrm{CH}}\nolimits}
\def\proetale{{pro\text{-}\acute{e}tale}}
\def\etale{{\acute{e}tale}}
\def\QCoh{\mathit{QCoh}}
\def\Ker{\mathop{\mathrm{Ker}}}
\def\Im{\mathop{\mathrm{Im}}}
\def\Coker{\mathop{\mathrm{Coker}}}
\def\Coim{\mathop{\mathrm{Coim}}}

% Boxtimes
%
\DeclareMathSymbol{\boxtimes}{\mathbin}{AMSa}{"02}

%
% Macros for moduli stacks/spaces
%
\def\QCohstack{\mathcal{QC}\!\mathit{oh}}
\def\Cohstack{\mathcal{C}\!\mathit{oh}}
\def\Spacesstack{\mathcal{S}\!\mathit{paces}}
\def\Quotfunctor{\mathrm{Quot}}
\def\Hilbfunctor{\mathrm{Hilb}}
\def\Curvesstack{\mathcal{C}\!\mathit{urves}}
\def\Polarizedstack{\mathcal{P}\!\mathit{olarized}}
\def\Complexesstack{\mathcal{C}\!\mathit{omplexes}}
% \Pic is the operator that assigns to X its picard group, usage \Pic(X)
% \Picardstack_{X/B} denotes the Picard stack of X over B
% \Picardfunctor_{X/B} denotes the Picard functor of X over B
\def\Pic{\mathop{\mathrm{Pic}}\nolimits}
\def\Picardstack{\mathcal{P}\!\mathit{ic}}
\def\Picardfunctor{\mathrm{Pic}}
\def\Deformationcategory{\mathcal{D}\!\mathit{ef}}
\usepackage{cleveref}
\crefname{originallemma}{original lemma}{original lemmas}
\Crefname{originallemma}{Original lemma}{Original lemmas}
\makeatletter
\newlabel{divisors-lemma-divisor-meromorphic-well-defined}{{02SH}{source}{}{sourcepointer-1}{}}
\newlabel{divisors-lemma-divisor-meromorphic-well-defined@cref}{{[originallemma][0][]02SH}{[1][0][]source}{}{}{}}
\makeatother
\crefname{originallemma}{original lemma}{original lemmas}
\Crefname{originallemma}{Original lemma}{Original lemmas}
\makeatletter
\newlabel{item-normalization}{{0EAN}{source}{}{sourcepointer-2}{}}
\newlabel{item-normalization@cref}{{[originallemma][0][]0EAN}{[1][0][]source}{}{}{}}
\makeatother
\crefname{originallemma}{original lemma}{original lemmas}
\Crefname{originallemma}{Original lemma}{Original lemmas}
\makeatletter
\newlabel{lemma-milnor-gersten-low-degree}{{0EAX}{source}{}{sourcepointer-3}{}}
\newlabel{lemma-milnor-gersten-low-degree@cref}{{[originallemma][0][]0EAX}{[1][0][]source}{}{}{}}
\makeatother
\title{The key formula identifying a difference of divisors weighted by orders with the divisor of the tame symbol}
\author{The Stacks Project Authors}
\date{stacks-project:chow.tex:local-snapshot-20261012; GNU FDL 1.2 or later}
\begin{document}
\maketitle
\noindent\textit{Editorial note.} All delivered mathematical text is the author's own, copied from the pinned Stacks Project chapter file chow.tex; no statement or proof step was rewritten, reordered or removed.\par\medskip
\noindent\textit{Editorial note.} Preamble provenance: the pinned chapter file opens with \textbackslash{}input\{preamble\} and carries no preamble text of its own, so the project's single Stacks-wide preamble staged at sources/stacks-chow/preamble.tex is excerpted verbatim as two spans, lines 14-18 (Xy-pic setup) and lines 154-217 (the native theorem declarations and mathematical macros). Nothing was added, removed or redefined; the wrapper class and the amsmath/amssymb/amsthm/graphicx base packages are supplied by the assembler.\par\medskip
\noindent\textit{Editorial note.} Same-chapter pointers are located at their \textbackslash{}label lines inside the pinned chapter file; cross-chapter pointers are located in the project's own tag list staged in the snapshot (sources/stacks-chow/tags.txt), at the verbatim lines carrying their canonical tags.\par\medskip
\noindent\textit{Editorial note.} The author's closing sentence at source line 4815, 'which is zero by the key Lemma \textbackslash{}ref\{lemma-milnor-gersten-low-degree\}.', is retained verbatim at the end of the counted proof; the reference is kept as a tagged navigation pointer. The coefficient comparison that produces the displayed sum is carried out inside the span, so only the final citation is left to the reader.\par\medskip
\noindent\textit{Editorial note.} Span note: the pool card recorded proof\_lines 4760-4815, one line short of the closing \textbackslash{}end\{proof\} at line 4816, which made the first assembly attempt fail with a runaway proof environment; the counted proof span is the real environment 4760-4816, taken from the pinned chapter file.\par\medskip
\noindent\textit{Allowed standard prerequisites (editorial).}\begin{itemize}
\item Well-definedness of the divisor of a meromorphic function. Divisors, label divisors-lemma-divisor-meromorphic-well-defined, invoked at proof line 4769. Justifies that replacing a trivialising section by a unit multiple changes the divisor by div(u|\_\{Z\_i\}), which is the computation in the first paragraph.
\item Normalisation of the tame symbol and its bilinearity. Same chapter chow.tex, property item-normalization of the tame symbol, invoked at proof line 4775, together with bilinearity used at proof lines 4772-4774. Shows the two extra terms produced by the replacement s\_i -> u s\_i are equal, which is what allows the reduction a\_i = b\_i = 1.
\item Vanishing of the sum of the orders of the tame symbol over the height one primes. Same chapter chow.tex, label lemma-milnor-gersten-low-degree, invoked at proof line 4815. Supplies the vanishing of the final sum of orders; the reduction of the coefficient comparison to that sum is carried out in the counted span.
\end{itemize}\medskip
\ifdefined\hypertarget\hypertarget{sourcepointer-1}{}\fi
\noindent\textit{Source reference.} Original lemma 02SH: sources/stacks-chow/tags.txt, source lines 3660--3660. Navigation pointer for \textbackslash{}ref\{divisors-lemma-divisor-meromorphic-well-defined\} used inside the counted span (line 4769); the argument that consumes it is retained in the span.\par\medskip
\ifdefined\hypertarget\hypertarget{sourcepointer-2}{}\fi
\noindent\textit{Source reference.} Original lemma 0EAN: sources/stacks-chow/pinned-original-source.tex, source lines 790--790. Navigation pointer for \textbackslash{}ref\{item-normalization\} used inside the counted span (line 4775); the argument that consumes it is retained in the span.\par\medskip
\ifdefined\hypertarget\hypertarget{sourcepointer-3}{}\fi
\noindent\textit{Source reference.} Original lemma 0EAX: sources/stacks-chow/pinned-original-source.tex, source lines 1263--1263. Navigation pointer for \textbackslash{}ref\{lemma-milnor-gersten-low-degree\} used inside the counted span (line 4815); the argument that consumes it is retained in the span.\par\medskip
\begin{lemma}[Key formula]
\label{lemma-key-formula}
In the situation above the cycle
$$
\sum
(Z_i \to X)_*\left(
\text{ord}_{B_i}(f_i) \text{div}_{\mathcal{N}|_{Z_i}}(t_i|_{Z_i}) -
\text{ord}_{B_i}(g_i) \text{div}_{\mathcal{L}|_{Z_i}}(s_i|_{Z_i}) \right)
$$
is equal to the cycle
$$
\sum (Z_i \to X)_*\text{div}(\partial_{B_i}(f_i, g_i))
$$
\end{lemma}

\begin{proof}
First, let us examine what happens if we replace $s_i$ by $us_i$
for some unit $u$ in $B_i$. Then $f_i$ gets replaced by $u^{-1} f_i$.
Thus the first part of the first expression of the lemma is unchanged
and in the second part we add
$$
-\text{ord}_{B_i}(g_i)\text{div}(u|_{Z_i})
$$
(where $u|_{Z_i}$ is the image of $u$ in the residue field) by
Divisors, Lemma \ref{divisors-lemma-divisor-meromorphic-well-defined}
and in the second expression we add
$$
\text{div}(\partial_{B_i}(u^{-1}, g_i))
$$
by bi-linearity of the tame symbol. These terms agree by property
(\ref{item-normalization}) of the tame symbol.

\medskip\noindent
Let $Z \subset X$ be an irreducible closed with $\dim_\delta(Z) = n - 2$.
To show that the coefficients of $Z$ of the two cycles of the lemma
is the same, we may do a replacement $s_i \mapsto us_i$ as in the previous
paragraph. In exactly the same way one shows that we may do a replacement
$t_i \mapsto vt_i$ for some unit $v$ of $B_i$.

\medskip\noindent
Since we are proving the equality of cycles we may argue one coefficient
at a time. Thus we choose an irreducible closed $Z \subset X$
with $\dim_\delta(Z) = n - 2$ and compare coefficients. Let $\xi \in Z$
be the generic point and set $A = \mathcal{O}_{X, \xi}$. This is a Noetherian
local domain of dimension $2$. Choose generators $\sigma$ and $\tau$
for $\mathcal{L}_\xi$ and $\mathcal{N}_\xi$. After shrinking $X$, we may
and do assume $\sigma$ and $\tau$ define trivializations
of the invertible sheaves $\mathcal{L}$ and $\mathcal{N}$ over all of $X$.
Because $Z_i$ is locally
finite after shrinking $X$ we may assume $Z \subset Z_i$ for all $i \in I$
and that $I$ is finite. Then $\xi_i$ corresponds to a prime
$\mathfrak q_i \subset A$ of height $1$.
We may write $s_i = a_i \sigma$ and $t_i = b_i \tau$
for some $a_i$ and $b_i$ units in $A_{\mathfrak q_i}$.
By the remarks above, it suffices to prove the lemma when
$a_i = b_i = 1$ for all $i$.

\medskip\noindent
Assume $a_i = b_i = 1$ for all $i$. Then the first expression of the
lemma is zero, because we choose $\sigma$ and $\tau$ to be trivializing
sections. Write $s = f\sigma$ and $t = g \tau$ with $f$ and $g$ in the
fraction field of $A$. By the previous paragraph we have reduced to the case
$f_i = f$ and $g_i = g$ for all $i$. Moreover, for a height $1$ prime
$\mathfrak q$ of $A$ which is not in $\{\mathfrak q_i\}$ we have
that both $f$ and $g$ are units in $A_\mathfrak q$ (by our choice of
the family $\{Z_i\}$ in the discussion preceding the lemma). Thus
the coefficient of $Z$ in the second expression of the lemma is
$$
\sum\nolimits_i \text{ord}_{A/\mathfrak q_i}(\partial_{B_i}(f, g))
$$
which is zero by the key Lemma \ref{lemma-milnor-gersten-low-degree}.
\end{proof}

\end{document}
